\documentclass[11pt]{article}

\usepackage[margin=1.1in]{geometry}
\usepackage[utf8]{inputenc}
\usepackage[T1]{fontenc}
\DeclareUnicodeCharacter{03B8}{\ensuremath{\theta}}
\DeclareUnicodeCharacter{2212}{\ensuremath{-}}
\usepackage{amsmath,amssymb,amsthm}
\usepackage{booktabs}
\usepackage{microtype}
\usepackage[hidelinks]{hyperref}

\newcommand{\E}{\mathbb{E}}
\newcommand{\Prob}{\mathbb{P}}
\newcommand{\code}[1]{\texttt{#1}}
\newtheorem{claim}{Claim}

\title{Smoothed Perturbation Analysis, From the Beginning\\
\large The derivation, the estimator we implement, and what re-declared
clocks change}
\author{Concourse research notes}
\date{July 17, 2026}

\begin{document}
\maketitle
\tableofcontents

%==============================================================================
\section{What this document is}

The design conversation asked for a derivation. SPA works in Concourse
today for models whose clocks keep one distribution for life, and it
refuses records that contain re-declared clocks, because the mathematics
for that case had not been worked out. This document is the first step of
working it out, following the process the project uses for such
questions: write the original derivation carefully, with notation defined
and with the reason for each step stated, and only then ask what changes
when clocks are re-declared. Books and articles on SPA are terse. This
report is deliberately not.

The document also folds in what the Question 3 experiments discovered
while it was being written. Those experiments were supposed to be a
side check on the pathwise estimator. Instead they located a structural
gap that sits underneath both the pathwise estimator and SPA, and the
multi-segment part of this report is organized around that finding.

Sections 2 through 4 build the single-segment theory from nothing.
Section 5 states the estimator exactly as implemented, so the theory and
the code can be compared line by line. Sections 6 and 7 treat re-declared
clocks: what generalizes cleanly, what breaks, a small closed-form result
that explains the experiments, and the ladder of falsification tests that
should gate any implementation. Section 8 collects the decisions.

%==============================================================================
\section{The setting and the notation}

\subsection{The race}

A generalized semi-Markov process runs as a race among clocks. The state
is $s$. In state $s$ a set of clocks is enabled. Clock $i$ was enabled at
time $\tau_i$ (its \emph{enabling time}) and carries a declared
probability distribution $F_i$ for its holding time --- the time from
$\tau_i$ until it would fire. The clock that fires first wins. Firing
moves the state, which enables new clocks and disables others, and the
race resumes. The model's entire statement about a clock is the pair
$(F_i, \tau_i)$, read as the conditional law of its firing time given the
history. Nothing else about the randomness is the model's business.

Parameters $\theta \in \mathbb{R}^D$ enter in exactly one place: the
declared distributions may depend on $\theta$, so we write
$F_i(\cdot\,;\theta)$. The performance measure is a functional $L$ of the
path --- a terminal reading, a time integral such as
$\int_0^T N(t)\,dt$, or a first-passage time. We want
$\partial_\theta \E[L]$.

\subsection{The path as a deterministic map}

Give clock $i$ a uniform random number $u_i$ and set its holding time by
inversion,
\[
X_i \;=\; F_i^{-1}(u_i\,;\theta),
\]
so that $X_i$ has law $F_i$. (When a clock is disabled and later
re-enabled, or fires repeatedly, it consumes a fresh uniform each time;
index them $u_{i,1}, u_{i,2}, \dots$. Nothing below depends on that
detail.) Once every uniform is fixed, the whole path is a deterministic
function of $\theta$: each firing time, each state visit, and hence
$L = L(\theta, u)$ where $u$ is the vector of all uniforms. The quantity
we want is
\[
\frac{d}{d\theta}\,\E[L] \;=\; \frac{d}{d\theta} \int L(\theta, u)\,du .
\]

Everything in perturbation analysis is a fight about whether the
derivative may be moved inside this integral.

\subsection{Where the map is smooth and where it is not}

Fix $u$ and move $\theta$ a little. Each holding time
$X_i = F_i^{-1}(u_i;\theta)$ slides smoothly, because the inverse
survival function of a fixed uniform is a smooth function of the
distribution's parameters. Firing times are sums and minima of holding
times, so they slide too. As long as the \emph{order} of firings does not
change, the path deforms continuously: the same events happen, at
slightly different times. On this region $L(\cdot, u)$ is typically
differentiable, and its derivative is what the pathwise (IPA) estimator
computes.

The map is not smooth everywhere. At certain values of $\theta$ two
firing times tie. Crossing such a value swaps the order of two events,
and the path after the swap can be a different path entirely: a different
job got the server, a different queue overflowed. If $L$ reads anything
that depends on the order --- a count, an indicator, a first passage ---
then $L(\cdot, u)$ jumps at the tie. The set of $u$ for which a tie
happens at a given $\theta$ has probability zero, which is why the
pathwise estimator is well defined; but the set of $u$ for which a tie
happens \emph{somewhere near} $\theta$ has probability of order
$\delta\theta$, and each such $u$ contributes a jump of order one. The
product is a term of order $\delta\theta$ that pure pathwise
differentiation cannot see:
\[
\frac{d}{d\theta}\,\E[L]
\;=\;
\underbrace{\E\!\left[\frac{\partial L}{\partial \theta}\right]}_{\text{IPA: the sliding of times}}
\;+\;
\underbrace{\text{(probability flux across ties)} \times \text{(jump size)}}_{\text{the boundary term IPA drops}} .
\]
The first term is exact when $L$ is continuous at order swaps (a time
integral of the state usually is). The second term is the whole story for
order-carried observables, and estimating it is what SPA is for.

%==============================================================================
\section{The smoothing idea}

\subsection{Conditioning makes the discontinuity differentiable}

The trick, due to Gong and Ho and developed at book length by Fu and Hu,
is not to differentiate $L$ itself but a conditional expectation of $L$.
Choose a \emph{characterization} $Z$: a set of random quantities that
pins down the path except for one distinguished piece of randomness.
Then
\[
\E[L] \;=\; \E\bigl[\,\E[L \mid Z]\,\bigr],
\]
and the inner quantity $\E[L \mid Z]$ is a function of $\theta$ that
integrates over the distinguished randomness. Integration smooths. Where
$L$ has a jump indicator $\mathbf{1}\{X_e < \xi\}$ in it, the conditional
expectation has the continuous function $\Prob(X_e < \xi \mid Z)$
instead. If the conditional expectation is differentiable in $\theta$ and
the interchange is justified, then
\[
\frac{d}{d\theta}\,\E[L]
\;=\;
\E\!\left[\frac{d}{d\theta}\,\E[L \mid Z]\right],
\]
and the inner derivative now contains the boundary contribution
explicitly, as the derivative of a probability. The estimator samples the
inner derivative. That is the entire idea. The art is in choosing $Z$ so
that the inner derivative is computable from one simulated path plus a
small amount of extra work.

\subsection{The characterization we use}

Concourse's SPA (following Fu and Hu's GSMP treatment) applies the idea
locally, at each firing of the simulation. Fix an epoch: the winner $e_k$
fires at time $t_k$. Consider one other clock $e$ that was enabled at
that moment --- a \emph{candidate} --- with enabling time $\tau_e$ and
age
\[
\xi \;=\; t_k - \tau_e
\]
at the decision. In the realized path the candidate lost: $X_e > \xi$.
In a nearby path it might have won, and that is the order swap whose
boundary contribution we must capture.

Condition on everything except the candidate's holding time $X_e$: all
other uniforms, hence all other holding times, hence (up to this epoch)
the winner's firing time $t_k$ and the candidate's enabling time
$\tau_e$. Under this conditioning, what do we know about $X_e$? We know
the candidate survived every earlier epoch --- at each previous firing it
was enabled and did not fire. All those survivals are implied by the
single statement $X_e > \xi$, because $\xi$ is its age at the latest
decision. So the conditional law of $X_e$ given the characterization is
its declared law conditioned on survival past $\xi$:
\[
\Prob(X_e \in dx \mid Z) \;=\; \frac{f_e(x;\theta)}{1 - F_e(\xi;\theta)}\,dx,
\qquad x > \xi .
\]
This is why the boundary weight will be a \emph{hazard} and not a
marginal density. The candidate is not a fresh clock; it is a clock we
have already watched survive to age $\xi$, and every probability
statement about it must respect that. The implementation pins this with
an exact two-clock-race check: with the hazard weight the boundary term
integrates to the true derivative $\lambda_b/(\lambda_a+\lambda_b)^2$ of
a win probability; with the marginal density it integrates to the wrong
value $\lambda_b/(2\lambda_a+\lambda_b)^2$.

%==============================================================================
\section{The single-segment derivation}

\subsection{The boundary probability and its derivative}

Work at one epoch with one candidate, everything else conditioned. Write
the conditional expectation of the functional as an integral over the
candidate's holding time $x$:
\[
\E[L \mid Z]
\;=\;
\int_{\xi(\theta)}^{\infty}
L(x;\theta)\,
\frac{f_e(x;\theta)}{1 - F_e(\xi(\theta);\theta)}
\,dx .
\]
Three things here depend on $\theta$, and keeping them separate is the
whole derivation. First, the integrand path $L(x;\theta)$: for $x$ away
from $\xi$ the order is fixed and $L$ deforms smoothly --- this produces
the IPA part, and we set it aside. Second, the density $f_e(\cdot;\theta)$
and its normalization: these produce terms that pair up with the IPA
part's likelihood bookkeeping and, in the implemented estimator, are
carried by the same replay. Third --- the term we came for --- the lower
limit $\xi(\theta)$, the boundary between ``candidate loses'' and
``candidate wins.''

Differentiating the integral at its moving lower limit gives the boundary
contribution
\[
-\,\frac{f_e(\xi;\theta)}{1 - F_e(\xi;\theta)}\;
\frac{d\xi}{d\theta}\;
L(\xi^+;\theta)
\;+\;
(\text{the matching term in which the boundary point } x = \xi
\text{ moves with the law}) .
\]
The second piece needs one more sentence of care. The boundary is a tie
between two \emph{times}, not a fixed number: the candidate wins when its
holding time, viewed at fixed uniform as $X_e(\theta)$, drops below the
age $\xi(\theta)$. Both sides move with $\theta$. The rate at which
probability crosses the boundary is proportional to the \emph{relative}
velocity
\[
v \;=\; \frac{d\xi}{d\theta} \;-\; \frac{dX_e}{d\theta}\Big|_{X_e = \xi},
\qquad
\frac{dX_e}{d\theta}\Big|_{X_e=\xi}
\;=\;
-\,\frac{\partial_\theta F_e(\xi;\theta)}{f_e(\xi;\theta)} ,
\]
where the last identity is the fixed-uniform (probability-integral
transform) derivative of $X_e = F_e^{-1}(u;\theta)$ evaluated at the
boundary point. When $v > 0$ the age is catching up to the candidate's
sample: order swaps flow \emph{into} the swapped configuration at rate
$h_e(\xi)\,v$, where
\[
h_e(\xi) \;=\; \frac{f_e(\xi;\theta)}{1 - F_e(\xi;\theta)}
\]
is the hazard. When $v \le 0$ the boundary recedes and no probability
crosses; this is why the implemented factor is $(v)^+$, the positive
part.

What multiplies the flux is the \emph{jump}: the difference between the
functional's value on the path where the pair fired in swapped order and
its value on the nominal path, in the limit where the holding time
between the two firings vanishes. Call these the perturbed path (PP:
candidate first) and the degenerated nominal path (DNP: winner first,
candidate immediately after, if still enabled). The local boundary
contribution at this epoch and candidate is
\[
\boxed{\;
h_e(\xi)\;\Bigl(\tfrac{d\xi}{d\theta} - \tfrac{dX_e}{d\theta}\big|_{\xi}\Bigr)^{\!+}\;
\bigl(L(\mathrm{PP}) - L(\mathrm{DNP})\bigr)
\;}
\]
and the SPA estimator is the IPA term plus the sum of these
contributions over epochs and candidates, plus one more family of terms
at the horizon (next subsection). Summing over epochs is legitimate
because each epoch's contribution came from conditioning on a different
characterization; Fu and Hu's Theorem 3.3 hypotheses (differentiable
conditional expectation, uniformly integrable difference quotients) are
what license both the interchange and the sum.

\subsection{Why the jump is estimated with a coupled clone pair}

$L(\mathrm{PP})$ and $L(\mathrm{DNP})$ are values of the functional on
two constructed continuations of the same prefix. The construction must
respect the characterization: everything except the candidate's residual
is held fixed, so both continuations start from the same state, fire the
same two events at the same instant $t_k$ (in opposite orders), and then
continue \emph{with the same future randomness}. In the implementation
both clones branch from the pre-commit world, receive one shared rekey
seed, force the pair in their respective orders, and run to the horizon.
Sharing the seed makes the difference $L(\mathrm{PP}) - L(\mathrm{DNP})$
quiet: the continuation randomness cancels and only the effect of the
order remains. An independent-seeds version would be unbiased too, but
its variance would swamp the term being estimated.

Two structural facts reduce the cost. First, the \emph{commuting gate}:
if firing the pair in both orders re-coalesces to the same state, with
both events surviving in both orders, then under the characterization the
two continuations are the same path and the jump is exactly zero --- no
clones needed. This is a correctness statement, not just a saving: the
clone pair in that case would return pure noise. (A first-passage
functional needs a stronger gate, because it reads the zero-duration
intermediate states.) Second, the \emph{horizon term}: a clock scheduled
just past the observation time $T$ crosses into the window as $\theta$
moves. That is a boundary with the constant $T$, so its ``decision-time
derivative'' is zero, no clones are needed, and the jump is read directly
off the fired state. A time integral has zero jump there (the crossing
event occupies zero width), consistent with IPA already being exact for
it.

\subsection{The velocity factor comes from the replay}

The term $d\xi/d\theta = dt_k/d\theta - d\tau_e/d\theta$ requires the
derivative of two realized \emph{times} with the order frozen --- exactly
what the pathwise replay computes. The implementation takes the Jacobian
of \code{replay\_times} at the recorded path and reads off the rows for
the winner's epoch and the candidate's enabling epoch. The term
$dX_e/d\theta|_\xi$ comes from the candidate's law directly
($-\partial_\theta F / f$ at the frozen age $\xi$), never from the
replay: the candidate did not fire in the realized path, so it has no
replayed time, and using a replayed dual age inside $\partial_\theta F$
would count $\theta$ twice.

This dependence is worth stating loudly, because Section 6 turns on it:
\emph{the boundary term is built on top of the pathwise replay}. If the
replay's time derivatives are wrong, the SPA weight multiplies wrong
velocities, and no boundary correction can rescue the sum.

\subsection{The validity conditions, collected}

For the single-segment estimator as implemented, the requirements are:

\begin{enumerate}
\item Each clock keeps one declared law from enabling to firing, so the
  conditional law given survival is read off that one law (the hazard at
  the age).
\item The clock laws are dual-safe: the replay's inverse survival
  functions differentiate under ForwardDiff.
\item \code{fire} is pure and states compare by value, so the commuting
  gate can fire speculatively.
\item After forcing the pair, each surviving clock's continuation law is
  its declared law conditioned on its own age --- the world's
  rekey-then-jitter performs exactly this resample-given-survival at a
  stopping time.
\item The functional and laws satisfy the interchange hypotheses
  (differentiable conditional expectations, uniformly integrable
  quotients). We do not verify these per model; the finite-difference
  comparisons in the charter are the operational check.
\end{enumerate}

%==============================================================================
\section{The implemented estimator, in the derivation's notation}

For reference, the map between the theory and \code{ClockGradients/src/spa.jl}:

\begin{center}
\small
\begin{tabular}{@{}ll@{}}
\toprule
Theory & Implementation \\
\midrule
IPA term & \code{ipa\_gradient} on the replication's own record \\
hazard $h_e(\xi)$ & \code{pdf(d,\,\char`\~xi)/ccdf(d,\,\char`\~xi)} at the enabling-state law \\
$d\xi/d\theta$ & rows of \code{ForwardDiff.jacobian(replay\_times)} \\
$dX_e/d\theta|_\xi$ & $-\,$\code{ForwardDiff.gradient(cdf)}$/f(\xi)$ \\
positive part $(v)^+$ & \code{max(dxi - dx, 0.0)} \\
jump $L(\mathrm{PP})-L(\mathrm{DNP})$ & coupled clone pair, shared rekey seed \\
commuting gate & \code{zero\_jump\_certified} (two speculative \code{fire}s) \\
horizon family & \code{\_horizon\_jump}, no clones \\
candidate enumeration & all enabled non-winners (\code{HazardWeight}), \\
& or next winner with truncated hazard (\code{TruncatedHazard}) \\
multi-segment refusal & \code{has\_chains(record)} throws \\
\bottomrule
\end{tabular}
\end{center}

The refusal in the last row is the subject of the rest of this document.
Per the design conversation, it stays a whole-record refusal: if any part
of a record is beyond the estimator's derivation, the user hears a clear
no rather than a quietly wrong number.

%==============================================================================
\section{What re-declared clocks change}

\subsection{The vocabulary}

Under processor sharing (and any future discipline that re-evaluates
laws mid-flight), a clock's declared distribution is replaced at event
times while the clock stays enabled. Each replacement starts a new
\emph{segment}. The enabling time $\tau_e$ stays fixed; what changes is
the declared law, which after a re-declaration says: given the work this
job has received (the banked age $a$) and the current sharing speed, here
is the conditional law of the firing age. In Concourse the declared
object is \code{SharedRemaining(base, bank, speed, shift)}, whose support
begins at the segment's start.

Write the segment boundaries of one clock as ages
$0 = a_0 < a_1 < \dots < a_{m}$ past $\tau_e$ (each $a_i$ is some other
firing's time minus $\tau_e$), with declared laws $G_1, \dots, G_{m+1}$
on the successive segments. The boundaries are event times of the same
race we are differentiating, so each $a_i$ depends on $\theta$.

\subsection{What generalizes without new mathematics}

\emph{The weight.} The characterization argument in Section 3.2 never
used the number of segments. Conditioning on everything except the
candidate's residual randomness, the candidate's conditional law at the
decision is its \emph{current} declared law given survival to its current
age. So the boundary weight for a re-declared candidate is the hazard of
the latest segment's declared law at the age since $\tau_e$:
\[
h_e(\xi) \;=\; \frac{g_{m+1}(\xi)}{1 - G_{m+1}(\xi)} ,
\]
with $G_{m+1}$ built at the state of the last re-declaration rather than
the state of the enabling. In the implementation this means the
candidate bookkeeping must track the last re-declaration state, a
mechanical change.

\emph{The clone construction.} By the declaration rule, the re-declared
laws after a forced firing depend only on the post-firing state and the
clocks' own histories. If the two orders of the forced pair re-coalesce
to the same state --- and for a processor-sharing station two completions
in either order leave the same job set --- then the survivors'
re-declared laws coincide in both arms, and the commuting gate's
conclusion (zero jump) still holds. Forcing \emph{always} re-declares
every survivor under PS, but it re-declares them identically in the two
arms whenever the states match. The gate survives.

\subsection{Where it actually breaks: the velocity, not the weight}

Section 4.3 said the boundary term is built on the pathwise replay. The
Question 3 experiments (Section 7) show that the pathwise replay is
biased on multi-segment records --- under both couplings, in different
directions, even when every segment declares the same conditional law.
So the multi-segment problem is not, first, an SPA problem. It is a
problem in the object SPA stands on.

The cleanest statement of what is missing comes from differentiating the
carry chain itself. The chain maps a would-be firing age across a
re-declaration by matching conditional survival: if under the old law
$G$ the clock would fire at age $\alpha$, and at boundary age $a$ the law
is replaced by $G^{\mathrm{new}}$, the new would-be age
$\alpha'$ solves
\[
\log S^{\mathrm{new}}(\alpha') \;=\;
\log S^{\mathrm{new}}(a) + \log S(\alpha) - \log S(a),
\qquad S = 1 - G .
\]
This is exactly the transformation the sampler applies at a re-enable,
and it is what \code{\_replay\_carry} evaluates. Now differentiate the
identity with respect to anything --- a component of $\theta$, with the
boundary age $a = a(\theta)$ moving too. Using
$\partial_x \log S = -h$ (the hazard) and writing dots for derivatives:
\[
\boxed{\;
\dot{\alpha}'
\;=\;
\frac{h^{\mathrm{new}}(a) - h(a)}{h^{\mathrm{new}}(\alpha')}\;\dot{a}
\;+\;
\frac{h(\alpha)}{h^{\mathrm{new}}(\alpha')}\;\dot{\alpha}
\;+\;
\frac{\partial_\theta \log S(\alpha) - \partial_\theta \log S(a)
      + \partial_\theta \log S^{\mathrm{new}}(a)
      - \partial_\theta \log S^{\mathrm{new}}(\alpha')}
     {h^{\mathrm{new}}(\alpha')}
\;}
\]
Read the first coefficient. The sensitivity of the \emph{boundary time}
enters the clock's replayed future weighted by the \emph{hazard jump}
across the re-declaration, $h^{\mathrm{new}}(a) - h(a)$. Three
consequences, and each one matches an experiment:

\begin{enumerate}
\item If a re-declaration does not change the conditional law
  ($h^{\mathrm{new}} \equiv h$ across the boundary), the coefficient is
  zero and the chain step is transparent: a degenerate multi-segment
  clock must replay exactly like a single-segment one.
\item If the re-declaration changes the speed (processor sharing), the
  coefficient is proportional to the speed change, and the boundary
  time's derivative $\dot a$ --- the replayed derivative of some other
  firing --- genuinely feeds the clock's future. Dropping it biases the
  replay.
\item The current implementation cannot express the first term. The
  segment laws are rebuilt from states folded at the \emph{recorded}
  Float64 times, so \code{bank}, \code{anchor}, and speed are frozen
  constants with no derivative; and the evaluation of
  $\log S^{\mathrm{new}}$ at the boundary lands exactly at the new law's
  support start, where the value is pinned to zero regardless of how the
  boundary moves. Half of the channel ($h(a)\,\dot a$, through the old
  law) flows; the other half ($h^{\mathrm{new}}(a)\,\dot a$, through the
  new law) is killed. The two halves are supposed to cancel exactly in
  the degenerate case --- killing one of them is why the degenerate
  experiment shows large per-path errors, not small ones.
\end{enumerate}

A one-line repair was tried and refuted: taking the formula branch at
the support boundary changes nothing, because round-off places the
evaluation on either side of the boundary and the frozen anchors remain
frozen either way. The repair that the algebra asks for is structural:
the segment parameters inside the replayed laws must be dual-valued
functions of the replayed times. Concretely, either the replay folds
states at dual times (so \code{bank} and \code{anchor} carry
derivatives), or the record stores each firing's segment boundary
\emph{indices} (which firing re-declared this clock) so the replay can
rebuild boundary ages as differences of replayed times --- a
representation question, and precisely Question 4 of the design brief,
now with a derivative-driven reason to answer it.

\subsection{The redraw coupling drops the channel entirely}

The same algebra explains the redraw coupling's numbers. Its recurrence
anchors each firing at the \emph{last} re-declaration:
$t = \tau_e + \Phi^{-1}(\log u + \Phi(t_{\mathrm{draw}} - \tau_e))$ with
$\Phi = \log S$ of the last segment's law. The conditioning term
$\Phi(t_{\mathrm{draw}} - \tau_e)$ evaluates at the last segment's own
start, where it is pinned to zero, so the derivative of
$t_{\mathrm{draw}}$ never enters: each firing's sensitivity reduces to
the last segment's residual only. On the M/M/1--PS experiment this
truncation shows up as $\partial_\mu \approx -10$ against a true
$-47.7$: the estimator sees only the tail of each service, not the
$\theta$-dependence of the whole shared history.

\subsection{The two-job worked example}

The design brief asked for a two-job hand example, and it is worth
recording because it shows where processor sharing hides its jumps.

Two jobs share one server, no further arrivals, service requirements
$S_1, S_2$ drawn from $F(\cdot;\theta)$. While both are present each is
served at speed $\tfrac12$; alone, at speed $1$. The first completion is
at $2\min(S_1,S_2)$; the second at $\min(S_1,S_2) + \max(S_1,S_2)$ ---
work conservation: the total busy period is $S_1 + S_2$ regardless of
order. The occupancy integral is
\[
\int_0^\infty N(t)\,dt
\;=\;
2 \cdot 2\min + \bigl(\min + \max - 2\min\bigr)
\;=\;
3\min(S_1,S_2) + \max(S_1,S_2),
\]
a continuous, symmetric, piecewise-smooth function of $(S_1, S_2)$. At
the near-tie $S_1 = S_2$ the swap exchanges two identical futures: the
jump $L(\mathrm{PP}) - L(\mathrm{DNP})$ is zero, and the commuting gate
would prove it (two completions in either order leave the empty
station). So for occupancy-type functionals of a plain PS station, the
boundary term contributes nothing at completion--completion ties, and
the entire derivative must come through the pathwise term --- through
$\min$ and $\max$ sliding with $\theta$, including the way each
completion's time depends on the \emph{other} job's completion through
the speed change. This is exactly the channel Section 6.3 shows the
replay to be missing. The lesson: for processor sharing the hard part of
SPA is not the exotic boundary weight; it is that the smooth part must
first be made correct. Jumps that need the weight arise one level up,
when the swap changes something downstream that does not re-coalesce ---
a routing decision, a blocking threshold, a class boundary.

%==============================================================================
\section{The experiments}

Three experiments, all with fixed seeds, all against common-random-number
central finite differences (FD) on the same seeds; details in
\code{notes/event\_loop.tex} and the pinned charter tests.

\subsection{M/M/1--PS, occupancy integral}

$\lambda = 1$, $\mu = 2$, horizon 50, 400 records. Every record carries
multi-segment chains.

\begin{center}
\small
\begin{tabular}{@{}lcc@{}}
\toprule
 & $\partial_\lambda$ & $\partial_\mu$ \\
\midrule
FD & $93.4 \pm 3.9$ & $-47.7 \pm 1.8$ \\
IPA, redraw & $57.7 \pm 3.4$ & $-10.0 \pm 0.1$ \\
IPA, carry & $206.0 \pm 11.8$ & $-36.2 \pm 1.2$ \\
\bottomrule
\end{tabular}
\end{center}

Both couplings biased, in different directions; duals flow with no
errors; the Float64 replay reproduces the recorded times. Pinned as
charter item F14.

\subsection{The degenerate-segment experiment}

The design conversation proposed the decisive control: multi-segment
records whose segments all declare the same conditional law. A
64-server PS station serves every job at speed 1 at this load
($\lambda = 3$, $\mu = 1$; a slowdown would need 65 simultaneous jobs),
so its law is identical to the same station under FCFS --- but the PS
records carry chains (300 of 300) while the FCFS records are
segment-free.

\begin{center}
\small
\begin{tabular}{@{}lcc@{}}
\toprule
$\partial_\mu$ (true $\approx -143$) & segment-free records & degenerate chains \\
\midrule
FD & $-143.1 \pm 1.0$ & (same law) \\
IPA, redraw & $-142.8 \pm 1.0$ & $-29.5 \pm 0.1$ \\
IPA, carry & $-142.8 \pm 1.0$ & $-100.4 \pm 0.6$ \\
\bottomrule
\end{tabular}
\end{center}

Identical dynamics, opposite verdicts. The bias lives entirely in the
record's segment representation, which is what Section 6.3's algebra
predicts: the degenerate case should be transparent
($h^{\mathrm{new}} \equiv h$), and it fails only because the
implementation expresses half of a cancelling pair. Pinned as the
sharpened F14 test.

\subsection{The refuted one-line repair}

Taking the formula branch at the support boundary of
\code{SharedRemaining} (strict instead of weak inequality) changed no
number in either experiment: bit-identical gradients. Diagnosis in
Section 6.3; kept here because a refuted repair is as informative as a
successful one, and because it rules out the cheap fix before anyone
tries it again.

%==============================================================================
\section{The falsification ladder}

The derivation dictates the order of implementation work, each step
gated by a test that can refute it.

\begin{enumerate}
\item \emph{Segment-aware replay.} Give the replay the missing channel:
  rebuild segment laws with dual boundary parameters (fold at dual
  times, or record segment boundary indices --- the Question 4
  decision). Gate: the degenerate-segment test flips from refutation to
  agreement, path by path --- carry on degenerate chains must equal the
  segment-free replay exactly, not just on average.
\item \emph{Pathwise on PS.} Gate: the M/M/1--PS experiment's carry
  column matches FD; then a Weibull-base PS variant (non-exponential,
  dual-safe) matches FD, which exercises the $\partial_\theta \log S$
  terms the exponential case partially hides.
\item \emph{SPA weight for re-declared candidates.} Implement the
  latest-segment hazard of Section 6.2, replacing the whole-record
  refusal only after step 2 holds. Gate: on a functional where PS
  boundary terms vanish (occupancy), SPA must agree with the repaired
  IPA; on an order-carried functional (terminal count through a
  downstream contended resource), SPA must match FD where plain IPA
  visibly does not.
\item \emph{Jump construction audit.} The clone pair under PS re-declares
  survivors in both arms; the commuting gate covers the re-coalescing
  cases. Gate: an adversarial model where the two orders do \emph{not}
  re-coalesce (a blocking threshold read between the pair), checked
  against FD.
\end{enumerate}

Until step 2 passes, both the SPA refusal and a matching refusal for
IPA on chained records are the user-facing truth. The design
conversation settled that refusals are whole-record: a record either is
within the derivation or is not.

%==============================================================================
\section{Decisions and open questions}

\begin{itemize}
\item The hazard weight generalizes to re-declared candidates as the
  latest declared law's hazard at the age since enabling, built at the
  last re-declaration state. No new mathematics; bookkeeping only.
\item The commuting gate survives re-declaration, by the declaration
  rule, whenever the forced pair re-coalesces.
\item The genuinely new mathematics is the chain derivative of Section
  6.3. Its hazard-jump coefficient identifies exactly which sensitivity
  the current replay cannot represent, predicts the degenerate
  experiment's failure, and reduces to the classical recurrence when no
  law changes.
\item The implementation decision it forces is representational ---
  Question 4's ``where does segment truth live'' --- and should be taken
  with this document in hand: the reconstruction keeps the record
  minimal (the filtration), so the preferred route is to make the
  reconstruction dual-capable (fold at dual times), and to revisit
  storing boundaries only if that proves too costly.
\item Open: the uniform-integrability hypotheses for the multi-segment
  interchange; whether the truncated-hazard strategy's residual
  ($\eta$) needs a segment-aware form; and $\theta$-dependent marks,
  which stay excluded until the mark derivative term is derived.
\end{itemize}

\end{document}
